\documentclass[11pt]{article}
\usepackage[margin=0.8in]{geometry}
\usepackage{amsmath,amssymb,amsthm,microtype}
\newtheorem{theorem}{Theorem}
\begin{document}
\begin{center}
{\Large A Growth Bound for Arbitrarily Interleaved Transfer Bridges}\par\medskip
Collatz proof project\quad---\quad September 5, 2026
\end{center}
Let $T$ be the shortcut Collatz map and $B(x)=9x+2$.
Consider a finite path starting at a positive integer $n$, using ordinary
$T$ steps and bridges $B$ in any order. Let $x$ be its endpoint, $t$ its
number of ordinary steps, $j$ its number of odd ordinary steps, and $b$
its number of bridges. Set $e=2j+4b$.
We impose no congruence or recursive-parameter restrictions on bridges.
This relaxation makes the exclusion below apply to every admissible
forward transfer-bridge path as well.

\begin{theorem}[Mixed-path growth bound]
Every such path satisfies
\[
 2^t x\leq 2^e n.
\]
\end{theorem}
\begin{proof}
An even step divides the current value by two. At a positive odd value $y$,
$2T(y)=3y+1\leq4y$, so an odd step increases size by at most a factor of two.
A bridge satisfies $9y+2\leq16y$ for $y\geq1$.
Multiplying these bounds gives
$x\leq2^{j-(t-j)+4b}n=2^{e-t}n$.
The Lean proof uses natural-number inequalities and induction on the path,
without negative powers, logarithms, or a convergence premise.
\end{proof}

\begin{theorem}[Cycle charge cannot decrease]
If the starting and ending values both lie in $\{1,2\}$, then
\[
 t+\mathbf1_{n=1}\leq e+\mathbf1_{x=1}.
\]
Equivalently, a symbolic path with initial slope exponent $p$ and initial
clock $s$ cannot lower
\[
 L(y,p,s)=2p-s+\mathbf1_{y=1}
\]
between its starting and ending cycle states.
\end{theorem}
\begin{proof}
Write $n=2^{1-\mathbf1_{n=1}}$ and
$x=2^{1-\mathbf1_{x=1}}$ in the growth bound and compare powers of two.
The symbolic slope exponent increases by $j+2b=e/2$ and the clock by $t$,
giving the equivalent charge comparison. Lean checks the natural-number
inequality by the four possible cycle endpoint cases.
\end{proof}

\paragraph{Consequence for transfer search.}
At clock $s$, suppose a source and target have cycle base values and their
symbolic endpoints are $x+3^p2^{D-s}Q$ and $y+3^q2^{D-s}Q$.
If $L(x,p,s)>L(y,q,s)$, no subsequent forward source path, using any number
of bridges, can match the target's ordinary continuation in both constant
and slope. The target remains in the cycle and preserves its charge.
At a fixed clock, equal symbolic endpoints require equal cycle state and
slope exponent, hence equal charge, contradicting the theorem.
This allows sound pruning without imposing a future bridge-count bound.

\paragraph{What changed from the earlier cycle experiment.}
Previously, the exclusion covered repeated $2\mapsto20$ excursions with
no bridges inside an excursion. The new theorem allows bridges anywhere
along a forward path, including inside excursions. The earlier finite
sample's 302 negative-gap base pairs therefore cannot be rescued by any
forward bridge path started only after both unmodified base orbits have
entered the cycle. Earlier auxiliary paths, inverse steps, and other
inference rules are outside this conclusion.

\paragraph{Verification and remaining goal.}
\begingroup\raggedright
The module \texttt{Collatz.BridgeGrowth} proves positivity, the growth bound,
and the cycle inequality. The search now removes the excluded states.
A complete comparison on parameters 0--1000 preserves certificate outcomes
and first-hit depths with pruning enabled or disabled; all eight affine
search tests pass. The census remains Python evidence. The selected Lean
axiom audit contains 217 declarations, using only standard foundations.
\par\endgroup
This result limits a proof-search mechanism; it does not establish global
transfer coverage or prove Collatz. The next useful extension must reach
these cases before the late-cycle obstruction applies or introduce a
transition not covered by this forward-path bound. No priority claim is made.
\end{document}
